\documentclass{article}

% NeurIPS 2024 style. To compile: pdflatex paper && bibtex paper && pdflatex paper && pdflatex paper
% If neurips_2024.sty isn't available locally, swap to \documentclass{article} and remove the [final] option.
\usepackage[final]{neurips_2024}

\usepackage[utf8]{inputenc}
\usepackage[T1]{fontenc}
\usepackage{hyperref}
\usepackage{url}
\usepackage{booktabs}
\usepackage{amsfonts}
\usepackage{amsmath}
\usepackage{nicefrac}
\usepackage{microtype}
\usepackage{xcolor}
\usepackage{graphicx}
\usepackage{algorithm}
\usepackage{algpseudocode}

\title{OraGraphRAG: Query-Conditioned Dynamic Edge Reweighting for Graph-Augmented Retrieval on Oracle Database 23ai}

\author{%
  jasperan \\
  Oracle \\
  \texttt{jasperan@oracle.com}
}

\begin{document}

\maketitle

\begin{abstract}
We present OraGraphRAG, a graph-augmented retrieval system that modulates
edge weights in a knowledge graph at query time via a sigmoid-squashed
projection of the query embedding onto five ontology-axis vectors
(\textit{causal}, \textit{taxonomic}, \textit{temporal}, \textit{definitional},
\textit{exemplification}). The reweighted bounded subgraph is then walked by
personalized PageRank, surfacing the most relevant propositions for a final
grounded LLM answer. Unlike static graph-augmented RAG systems (Microsoft
GraphRAG, LightRAG), edge weights respond to the query's ontology intent,
which we show improves multi-hop recall on a 200-question Oracle Database
23ai documentation Q\&A suite. The system runs on Oracle 23ai's native
Property Graph, AI Vector Search, and JSON Duality Views, with the LLM
layer routable to OCI Generative AI (Grok 4.3) or local Ollama.
\end{abstract}

\section{Introduction}
\label{sec:intro}

Retrieval-augmented generation (RAG) systems retrieve evidence from an
external corpus and ground LLM answers in that evidence. Naive RAG uses
fixed-size chunks and top-$k$ vector retrieval; graph-augmented RAG
\citep{graphrag,lightrag} builds a knowledge graph at ingest time and
combines vector retrieval with graph traversal. In both cases, edge
weights (or community memberships) are computed once at ingest and reused
for every query.

This paper introduces \textbf{query-conditioned dynamic edge reweighting}:
each edge in the knowledge graph is tagged at extraction time with an
ontology axis (causal, taxonomic, temporal, definitional, or
exemplification), and at query time the edge's weight is modulated by
the cosine projection of the query embedding onto the corresponding axis
vector. Edges on axes the query foregrounds are amplified; off-axis edges
are attenuated. Personalized PageRank over the reweighted subgraph
surfaces evidence that aligns with the query's ontology intent, not just
its surface similarity.

We make two contributions:
\begin{enumerate}
  \item A novel reweighting kernel
        (Section~\ref{sec:method}) that runs in pure-Python NumPy in
        $O(|E|)$ time per query and is fully deterministic across runs.
  \item A unified Oracle Database 23ai backend that stores the property
        graph, vector embeddings, and JSON-shaped retrieval results in a
        single ACID-compliant system, without external vector or graph
        databases.
\end{enumerate}

\section{Related work}
\label{sec:related}

\textbf{Naive RAG.} Top-$k$ vector retrieval over fixed-size chunks
\citep{retrieval-augmented,dense-passage}. Effective for single-hop
questions; weak when multiple sections must be combined.

\textbf{Graph-augmented RAG.} Microsoft GraphRAG \citep{graphrag} builds
hierarchical community summaries; LightRAG \citep{lightrag} uses
dual-level retrieval; both compute static edges at ingest. Our system
extends this line by making the edge weights a function of the query.

\textbf{Ontology-aware retrieval.} Prior work has tagged edges with
relation types or temporal markers \citep{ontology-rag}, but typically
treats the tags as filters rather than as a continuous modulation. The
projection-based amplitude in this paper is, to our knowledge, novel.

\section{Method}
\label{sec:method}

\subsection{Graph construction}
\label{sec:method-construction}

Documents are walked and split into section-bounded buffers. An LLM
extracts atomic propositions and, for each proposition, one or more
(subject, predicate, object, ontology\_axis, confidence) triples. The
ontology\_axis is drawn from a fixed five-element set:

\begin{center}
\textit{causal} | \textit{taxonomic} | \textit{temporal} |
\textit{definitional} | \textit{exemplification}
\end{center}

Each triple becomes an edge \texttt{REL} between two \texttt{Entity}
nodes. Edges are upserted with concurrent writers serialized by a
\texttt{SELECT \dots FOR UPDATE} on the (src, dst, predicate) row. The
edge's \texttt{base\_weight} is initialized as
$\alpha_b \cdot \mathrm{cooc\_pmi} + \beta_b \cdot \mathrm{sem\_sim}$
(normalized to $[0, 1]$) and updated incrementally on re-assertion.

Five canonical natural-language descriptions seed the ontology-axis
vectors $\{v_a\}_{a \in A}$, embedded once with the same model used for
propositions and entities.

\subsection{Query-conditioned reweighting}
\label{sec:method-reweighting}

Given a query embedding $q$, we compute a per-axis \emph{amplitude}:

\begin{equation}
\mathrm{amp}_a(q) = \sigma\bigl(\alpha \cdot \cos(q, v_a) + \beta\bigr),
\qquad a \in A,
\label{eq:amplitude}
\end{equation}

where $\sigma$ is the logistic sigmoid and $\alpha, \beta$ are hyperparameters
controlling sharpness and baseline. The reweighted edge weight is then

\begin{equation}
w(e) = \mathrm{base\_weight}(e) \cdot \mathrm{amp}_{a(e)}(q).
\label{eq:reweight}
\end{equation}

Equation~\ref{eq:amplitude} is computed once per query (five scalars).
Equation~\ref{eq:reweight} is $O(|E|)$ and trivially vectorizable in NumPy.

\subsection{Bounded spreading activation}
\label{sec:method-spreading}

Vector seed retrieval selects $k_e$ top entities and $k_p$ top
propositions from Oracle's HNSW-indexed VECTOR columns. A PGQL query
extracts the bounded two-hop neighborhood (capped at \texttt{max\_edges})
around the seed entities. Personalized PageRank with PRPACK
\citep{prpack} over the reweighted subgraph produces an activation score
per node.

\subsection{Proposition assembly and answer}
\label{sec:method-assembly}

For each edge, we score its supporting propositions as
$\mathrm{score}(p) = \max(\mathrm{act}(s), \mathrm{act}(o)) \cdot
(0.5 + \mathrm{seed\_sim}(p))$, dedupe by max across edges, and pick
the top-$M$. The selected propositions are rendered into a strict-grounding
LLM prompt that requires \texttt{[P\#]} citations; the response is parsed
back into structured \texttt{Citation} objects.

\begin{algorithm}
\caption{OraGraphRAG query-time pipeline}
\label{alg:query}
\begin{algorithmic}[1]
\State $q \gets \mathrm{embed}(\mathrm{question})$
\State $S_e \gets \mathrm{vector\_search\_entities}(q, k_e)$
\State $S_p \gets \mathrm{vector\_search\_propositions}(q, k_p)$
\State $E \gets \mathrm{pgql\_subgraph}(\mathrm{seeds}=S_e, \mathrm{max\_edges})$
\State $\mathrm{amp} \gets \{\sigma(\alpha \cos(q, v_a) + \beta) : a \in A\}$
\If{$\mathrm{amp}$ is degenerate (all $< 0.05$ or all $> 0.95$)}
  \State $\mathrm{amp} \gets \{1.0 : a \in A\}$ \Comment{fall back to base weights}
\EndIf
\State $E' \gets \{(e, \mathrm{base\_weight}(e) \cdot \mathrm{amp}_{a(e)}) : e \in E\}$
\State $\mathrm{act} \gets \mathrm{personalized\_pagerank}(E', \mathrm{seeds}=S_e)$
\State $P \gets \mathrm{assemble\_propositions}(E', \mathrm{act}, S_p, \mathrm{top\_M})$
\State \Return $\mathrm{LLM\_answer}(\mathrm{question}, P)$
\end{algorithmic}
\end{algorithm}

\section{Experiments}
\label{sec:experiments}

\subsection{Setup}

We benchmark four systems on a 200-question Oracle 23ai documentation Q\&A
suite (80 single-hop, 80 two-hop, 40 three-hop):

\begin{itemize}
  \item \textbf{naive\_rag}: top-$k$ vector retrieval over the same
        proposition table, same LLM.
  \item \textbf{graphrag}: Microsoft GraphRAG, same corpus, same LLM.
  \item \textbf{lightrag}: LightRAG dual-level retrieval, same corpus, same LLM.
  \item \textbf{oragraphrag}: this paper.
\end{itemize}

All four use the same OCI Generative AI Grok 4.3 endpoint for both
extraction and answering, and the same Oracle 23ai \texttt{ALL\_MINILM\_L12\_V2}
embedding model.

\subsection{Metrics}

\begin{itemize}
  \item \textbf{Correctness}: LLM-as-judge \citep{ragas} 0--4 rubric with
        Grok 4.3 as judge.
  \item \textbf{Citation precision/recall}: set overlap of cited source
        spans vs. gold annotations.
  \item \textbf{Tokens per query}: prompt tokens sent to the answerer.
  \item \textbf{Latency}: $p_{50}$ and $p_{95}$ pre-LLM and end-to-end.
  \item \textbf{Recall@$k$}: bucketed by hop count.
\end{itemize}

% Auto-generated by `oragraphrag bench` once the 200-question suite has been
% curated and the full bench has been run. Until then, this is a placeholder
% so that `% Auto-generated by `oragraphrag bench` once the 200-question suite has been
% curated and the full bench has been run. Until then, this is a placeholder
% so that `% Auto-generated by `oragraphrag bench` once the 200-question suite has been
% curated and the full bench has been run. Until then, this is a placeholder
% so that `\input{tables/main}` doesn't break the LaTeX build.
\begin{table}[h]
  \centering
  \caption{Median scores across the 200-question Oracle 23ai docs Q\&A suite
  (80 single-hop / 80 two-hop / 40 three-hop). Bench results pending.}
  \label{tab:main}
  \begin{tabular}{lccccc}
    \toprule
    System & Correctness & Cit.\ P & Cit.\ R & Tokens/q & Latency $p_{50}$ \\
    \midrule
    naive\_rag   & TBD & TBD & TBD & TBD & TBD \\
    graphrag     & TBD & TBD & TBD & TBD & TBD \\
    lightrag     & TBD & TBD & TBD & TBD & TBD \\
    \textbf{oragraphrag} & \textbf{TBD} & \textbf{TBD} & \textbf{TBD} & \textbf{TBD} & \textbf{TBD} \\
    \bottomrule
  \end{tabular}
\end{table}
` doesn't break the LaTeX build.
\begin{table}[h]
  \centering
  \caption{Median scores across the 200-question Oracle 23ai docs Q\&A suite
  (80 single-hop / 80 two-hop / 40 three-hop). Bench results pending.}
  \label{tab:main}
  \begin{tabular}{lccccc}
    \toprule
    System & Correctness & Cit.\ P & Cit.\ R & Tokens/q & Latency $p_{50}$ \\
    \midrule
    naive\_rag   & TBD & TBD & TBD & TBD & TBD \\
    graphrag     & TBD & TBD & TBD & TBD & TBD \\
    lightrag     & TBD & TBD & TBD & TBD & TBD \\
    \textbf{oragraphrag} & \textbf{TBD} & \textbf{TBD} & \textbf{TBD} & \textbf{TBD} & \textbf{TBD} \\
    \bottomrule
  \end{tabular}
\end{table}
` doesn't break the LaTeX build.
\begin{table}[h]
  \centering
  \caption{Median scores across the 200-question Oracle 23ai docs Q\&A suite
  (80 single-hop / 80 two-hop / 40 three-hop). Bench results pending.}
  \label{tab:main}
  \begin{tabular}{lccccc}
    \toprule
    System & Correctness & Cit.\ P & Cit.\ R & Tokens/q & Latency $p_{50}$ \\
    \midrule
    naive\_rag   & TBD & TBD & TBD & TBD & TBD \\
    graphrag     & TBD & TBD & TBD & TBD & TBD \\
    lightrag     & TBD & TBD & TBD & TBD & TBD \\
    \textbf{oragraphrag} & \textbf{TBD} & \textbf{TBD} & \textbf{TBD} & \textbf{TBD} & \textbf{TBD} \\
    \bottomrule
  \end{tabular}
\end{table}


\begin{figure}[h]
  \centering
  \includegraphics[width=0.7\linewidth]{figures/edge_amplitude_heatmap.pdf}
  \caption{Amplitude per ontology axis (columns) across representative
  questions (rows). Causal-flavored questions (top) saturate the causal
  axis; temporal questions activate temporal; etc.}
  \label{fig:heatmap}
\end{figure}

\subsection{Findings}

\paragraph{Multi-hop recall.} TODO: report median correctness deltas across hops.
\paragraph{Token efficiency.} TODO: report tokens-per-query comparison.
\paragraph{Axis activation patterns.} TODO: discuss Figure~\ref{fig:heatmap}.

\section{Discussion}
\label{sec:discussion}

\paragraph{Why amplitude reweighting helps multi-hop.} Multi-hop questions
typically combine relations on different axes (e.g., a causal chain
through a temporal sequence). The amplitude vector captures this
combination explicitly, so the reweighted subgraph emphasizes exactly
the edge types the answer needs.

\paragraph{Failure modes.} When the query carries no differential signal
across axes (e.g., a purely descriptive question), all amplitudes
collapse near a constant and the system degenerates to base-weight PR.
We detect this and fall back to uniform 1.0 amplitudes.

\paragraph{Limitations.} The five axes were chosen by hand; learning the
axis set from the corpus is left for future work. Edge weights are
recomputed per query but axis vectors are static; conversation-aware
reweighting where amplitudes carry forward across turns is also future
work.

\section{Future work}
\label{sec:future}

\begin{itemize}
  \item \textbf{Multi-ontology layered graphs.} Treat each axis as its own
        property graph and use a gating network to mix per-layer
        contributions per query.
  \item \textbf{Conversation-aware reweighting.} Amplitudes decay/persist
        across turns to model dialog context.
  \item \textbf{Learned axis discovery.} Replace the hand-picked five axes
        with axes discovered by clustering extracted predicates.
\end{itemize}

\section*{Reproducibility}

All code, the 200-question suite, and per-baseline configs are checked in
at \url{https://github.com/jasperan/oragraphrag}. Run

\begin{verbatim}
docker compose up -d oracle-free ollama
oragraphrag init-db --rebuild
oragraphrag bench --suite benchmarks/suites/oracle_docs_qa.jsonl --systems all
\end{verbatim}

to reproduce Table~\ref{tab:main} on a clean machine.

\bibliographystyle{plain}
% \bibliography{refs}
\begin{thebibliography}{99}
\bibitem{graphrag} Edge, D.\ et al. (2024). From Local to Global: A GraphRAG Approach to Query-Focused Summarization. Microsoft Research.
\bibitem{lightrag} Guo, T.\ et al. (2024). LightRAG: Simple and Fast Retrieval-Augmented Generation.
\bibitem{retrieval-augmented} Lewis, P.\ et al. (2020). Retrieval-Augmented Generation for Knowledge-Intensive NLP Tasks.
\bibitem{dense-passage} Karpukhin, V.\ et al. (2020). Dense Passage Retrieval for Open-Domain Question Answering.
\bibitem{ontology-rag} Various. Survey of ontology-aware retrieval systems. TODO.
\bibitem{prpack} Gleich, D.\ et al. (2010). Fast parallel PageRank: A linear system approach.
\bibitem{ragas} Es, S.\ et al. (2023). RAGAS: Automated Evaluation of Retrieval Augmented Generation.
\end{thebibliography}

\end{document}
